\chapter{줄 세우기 비교}

엔도스랭크와 송금 랭크로 지갑 5,521개에 점수를 매겼어요.
이제 궁금한 것이 생겨요.
두 점수는 얼마나 비슷할까요?
비슷하다면 굳이 새 방법을 만들 필요가 없겠지요.

\section{두 줄이 얼마나 닮았나}

체육 시간에 다섯 친구가 줄을 섰어요.
한 번은 키 큰 순서로, 한 번은 나이 많은 순서로요.

\begin{figure}[h]
\centering
\begin{tikzpicture}[x=1.9cm]
  \node[font=\small, anchor=east] at (-0.2,1.6) {키 순서};
  \node[font=\small, anchor=east] at (-0.2,0) {나이 순서};
  \foreach \i/\k/\n in {0/a/지우, 1/b/하준, 2/c/서연, 3/d/도윤, 4/e/민서}
    \node[kid, minimum size=9mm, font=\scriptsize] (t\k) at (\i,1.6) {\n};
  \foreach \i/\k/\n in {0/b/하준, 1/a/지우, 2/c/서연, 3/e/민서, 4/d/도윤}
    \node[kid, minimum size=9mm, font=\scriptsize] (b\k) at (\i,0) {\n};
  \draw[line width=1pt, color=leafline] (tc) -- (bc);
  \foreach \k in {a,b,d,e}
    \draw[line width=1pt, color=roseline] (t\k) -- (b\k);
\end{tikzpicture}
\caption{친구 한 명마다 두 줄에서 선 자리를 선으로 이었어요. 선은 친구 수만큼 다섯 개이고, 엇갈린 곳이 두 군데 있어요.}
\end{figure}

두 줄을 비교할 때는 친구를 한 명씩 보지 않고, \emph{두 명씩 짝}을 지어서 봐요.
그림의 선 다섯 개는 친구 한 명에 하나씩이에요.
짝 하나는 친구 두 명, 그러니까 선 두 개를 고른 거예요.

다섯 명에서 두 명을 고르는 방법은 몇 가지일까요?
지우는 나머지 네 명과 짝이 될 수 있으니 4쌍이에요.
하준이는 지우와의 짝을 이미 셌으니 나머지 세 명과 3쌍,
서연이는 2쌍, 도윤이는 민서와 1쌍이에요.
모두 $4+3+2+1=10$, 열 쌍이에요.

쌍마다 이렇게 물어요.
“키 순서에서 앞에 선 친구가 나이 순서에서도 앞에 섰나?”

\begin{itemize}
\item 그렇다면 \textbf{맞는 쌍}이에요. 지우와 서연이는 두 줄 모두 지우가 앞이니 맞는 쌍이에요.
\item 아니라면 \textbf{어긋난 쌍}이에요. 지우와 하준이는 키로는 지우가 앞인데 나이로는 하준이가 앞이에요.
\end{itemize}

열 쌍을 모두 따져 보면 이래요.

\begin{center}
\small
\begin{tabular}{@{}ll@{\qquad}ll@{}}
\toprule
쌍 & 결과 & 쌍 & 결과 \\
\midrule
지우, 하준 & \textcolor{roseline}{어긋난 쌍} & 하준, 도윤 & 맞는 쌍 \\
지우, 서연 & 맞는 쌍 & 하준, 민서 & 맞는 쌍 \\
지우, 도윤 & 맞는 쌍 & 서연, 도윤 & 맞는 쌍 \\
지우, 민서 & 맞는 쌍 & 서연, 민서 & 맞는 쌍 \\
하준, 서연 & 맞는 쌍 & 도윤, 민서 & \textcolor{roseline}{어긋난 쌍} \\
\bottomrule
\end{tabular}
\end{center}

어긋난 쌍은 (지우, 하준)과 (도윤, 민서) 둘뿐이고, 나머지 여덟 쌍은 맞아요.
그림에서 두 친구의 선이 서로 엇갈려 있으면 그 두 친구는 어긋난 쌍이에요.
엇갈린 곳이 두 군데이니 어긋난 쌍도 둘이에요.

이제 점수를 만들어요.
\[
\text{닮은 정도} = \frac{\text{맞는 쌍} - \text{어긋난 쌍}}{\text{모든 쌍}} = \frac{8-2}{10} = 0.6
\]

이 점수를 \textbf{켄달의 타우}라고 부르고, 그리스 문자 $\tau$로 써요.
1938년 영국의 통계학자 모리스 켄달이 만들었어요.
두 줄이 완전히 같으면 1, 완전히 거꾸로면 $-1$, 아무 관계가 없으면 0 근처가 나와요.
이 책에서 앞으로 나오는 비교는 거의 모두 이 $\tau$로 해요.

\section{엔도스랭크와 송금 랭크는 얼마나 닮았을까?}

지갑 5,521개를 엔도스랭크 점수 순서로 한 줄, 송금 랭크 점수 순서로 한 줄 세웠어요.
두 줄의 $\tau$는 \textbf{0.360}이었어요.

1이 아니에요.
0.6도 아니에요.
두 점수는 조금 닮았지만, 지갑을 꽤 다르게 줄 세운다는 뜻이에요.
같은 지갑들을 보고도 “누가 믿을 만한가”에 대해 다른 대답을 하는 거예요.
이것이 논문의 첫 번째 질문에 대한 답이에요.

\begin{deeper}[0.360을 조금 더 들여다보기]
\setlength{\parskip}{0.55em}%
엔도스랭크 점수로 지갑 5,521개를 줄 세우면, 줄이 세 무리로 나뉘어요.

\begin{itemize}
\item \textbf{가 무리(131개)}: 연구 기간에 누군가에게서 허락을 \emph{받은} 지갑이에요.
      들어오는 화살표로 쪽지가 흘러 들어오니 점수가 아래의 기본 점수보다 높아요.
\item \textbf{나 무리(4,424개)}: 허락을 \emph{해 주기만} 하고 받은 적은 없는 지갑이에요.
      허락 지도에 점으로는 있지만 들어오는 화살표가 없어요.
      6장의 지우처럼 선생님께 받는 몫뿐이라서, 4,424개가 모두 똑같은 점수를 받아요.
      이 점수를 \emph{기본 점수}라고 부를게요.
\item \textbf{다 무리(966개)}: 허락 지도에 점으로도 나타나지 않는 지갑이에요.
      802개는 연구 기간에 허락을 하지도 받지도 않았고, 164개는 허락을 했지만 연구 기간이 끝날 때 그 허락이 모두 0이 되어 화살표가 지워졌어요(8장).
      논문의 주된 분석은 이렇게 지도 밖에 있는 지갑에 점수 0을 주었어요.
\end{itemize}

켄달의 $\tau$를 셀 때 규칙이 하나 더 있어요.
한 줄에서 두 지갑의 점수가 똑같으면, 그 쌍은 맞는 쌍도 어긋난 쌍도 아닌 \textbf{비긴 쌍}이라서 세지 않아요.
논문에서 쓴 $\tau_b$가 이렇게 세요.
그래서 엔도스랭크 줄에서 앞뒤가 정해지는 쌍은 서로 다른 무리에 속한 두 지갑의 쌍과, 가 무리 안에서 점수가 다른 쌍뿐이에요.

그 가운데 가장 많은 것이 나 무리와 다 무리 사이의 쌍이에요.
$4{,}424\times966$, 약 427만 쌍이에요.
엔도스랭크는 이 쌍에서 늘 나 무리 지갑을 앞에 세워요.
송금 랭크도 이 쌍의 88\%에서 같은 순서였어요.
허락 지도 밖의 지갑은 송금도 적게 주고받는 편이라 송금 랭크에서도 뒤쪽에 서거든요.
다 무리의 38\%는 송금 랭크 점수가 아예 0이에요.
0.360의 닮음은 사실상 모두 이 쌍들에서 나와요.
나머지 쌍들은 모두 합치면 오히려 어긋난 쌍이 조금 더 많아요.
다시 말해 두 점수는 “허락 지도에 있는 지갑인가, 없는 지갑인가”에서 같은 대답을 한 거예요.

그런데 지도 밖의 지갑에 0점을 주는 것은 여러 방법 가운데 하나를 고른 것이에요.
그래서 논문에서는 규칙을 바꾸어 다시 계산해 보았어요.
지도 밖의 966개 지갑도 허락 지도에 점으로 그려 넣되, 화살표는 하나도 달지 않는 거예요.
그리고 똑같이 쪽지 돌리기 놀이를 해요.
선생님은 지도에 있는 모든 점에 쪽지를 고르게 나누어 주니까, 이 점들도 이제 0이 아니라 선생님 몫을 받아요.
들어오는 화살표가 없으니 받는 것은 그것뿐이고, 그래서 나 무리와 똑같은 기본 점수가 돼요.
다 무리가 나 무리에 합쳐져서 5,390개 지갑이 모두 같은 점수가 되는 거예요.

\begin{center}
\begin{tikzpicture}[font=\scriptsize]
  \def\W{8.4}
  \def\ga{131/5521*\W}
  \def\na{4555/5521*\W}
  \node[anchor=south west] at (0,0.22) {주된 분석: 지도 밖 지갑은 0점};
  \fill[allowcol!85] (0,-0.2) rectangle ({\ga},0.2);
  \fill[allowcol!25] ({\ga},-0.2) rectangle ({\na},0.2);
  \fill[greycol!35] ({\na},-0.2) rectangle (\W,0.2);
  \node at ({(\ga+\na)/2},0) {나 4,424개: 모두 기본 점수};
  \node at ({(\na+\W)/2},0) {다 966개: 0점};
  \node[anchor=south west] at (0,-0.83) {바꾼 규칙: 지도 밖 지갑도 화살표 없는 점으로};
  \fill[allowcol!85] (0,-1.25) rectangle ({\ga},-0.85);
  \fill[allowcol!25] ({\ga},-1.25) rectangle (\W,-0.85);
  \node at ({(\ga+\W)/2},-1.05) {나 + 다 5,390개: 모두 기본 점수};
  \draw ({\ga/2},-1.25) -- ({\ga/2},-1.45) node[anchor=north west, xshift=-1mm] {가 131개: 기본 점수보다 높음 (줄의 맨 앞)};
\end{tikzpicture}
\end{center}

이렇게 하면 427만 쌍은 모두 비긴 쌍이 되어 빠져요.
남는 것은 가 무리 131개가 들어간 쌍, 약 71만 쌍뿐이에요.
이 쌍에서 송금 랭크가 엔도스랭크와 같은 순서인 것은 44\%, 반대 순서인 것은 56\%였어요.
허락을 받은 131개 지갑이 송금 랭크에서는 평균으로 나 무리보다 오히려 조금 뒤에 섰거든요.
맞는 쌍보다 어긋난 쌍이 조금 많으니 $\tau$는 $-0.025$, 거의 0이 돼요.

정리하면, 0.360은 엔도스랭크가 쪽지 놀이로 매긴 순서가 송금 랭크와 닮았다는 뜻이라기보다, 허락 지도 밖의 지갑을 두 점수가 모두 뒤쪽에 세웠다는 뜻에 가까워요.
엔도스랭크가 기본 점수보다 높게 구별해 낸 131개 지갑에서는 두 점수가 거의 관계가 없어요.

송금 랭크 쪽은 이 규칙으로 바뀌지 않아요.
송금 랭크에서는 선생님이 송금을 보낸 적 있는 지갑에게만 쪽지를 주기 때문에, 송금 지도 밖의 지갑 381개를 화살표 없는 점으로 넣어도 쪽지를 하나도 받지 못해 점수가 0 그대로예요.

두 규칙 가운데 어느 쪽이 옳다고 할 수는 없어요.
같은 점수를 어떻게 다루느냐에 따라 $\tau$가 이만큼 달라진다는 것을 알리려고 두 결과를 함께 실었어요.
바꾼 규칙은 결과를 본 뒤에 시도한 것이라 참고로만 보고해요.
\end{deeper}

\section{점수마다 무엇과 닮았나}

두 점수가 다르다는 건 알았어요.
그럼 각 점수는 무엇을 잘 나타낼까요?
이것을 알아보려고 지갑마다 아주 단순한 사실들을 세었어요.
이런 사실들을 \textbf{검증 지표}라고 불러요.

\begin{itemize}
\item \textcolor{sendcol}{\textbf{송금 검사}}: 송금을 몇 번 받았는지, 얼마나 받았는지
\item \textcolor{allowcol}{\textbf{허락 검사}}: 허락을 몇 번 받았는지, 얼마나 받았는지
\end{itemize}

이번에는 엔도스랭크와 송금 랭크만 보지 않고, 9장에서 만든 점수까지 모두 한 줄에 세웠어요.
그리고 점수 줄마다 검사 줄과의 $\tau$를 쟀어요.
결과는 이랬어요.

\begin{figure}[h]
\centering
\begin{tikzpicture}
\begin{axis}[
  xbar, bar width=6pt, width=0.92\linewidth, height=7.4cm,
  xmin=-0.13, xmax=0.78, xtick={0,0.2,0.4,0.6}, xlabel={$\tau$},
  symbolic y coords={A, L0, S, C, L1, E}, ytick={A, L0, S, C, L1, E},
  yticklabels={송금 랭크, 뒷면 랭크, 씨앗 랭크, 동전 랭크, 앞면 랭크, 엔도스랭크},
  enlarge y limits=0.11,
  extra x ticks={0}, extra x tick labels={}, extra x tick style={grid=major, grid style={greycol!70}},
  legend style={font=\scriptsize, at={(0.5,1.02)}, anchor=south, legend columns=2, draw=none},
  nodes near coords, nodes near coords style={font=\tiny},
  every node near coord/.append style={/pgf/number format/fixed, /pgf/number format/fixed zerofill, /pgf/number format/precision=3},
  tick label style={font=\small}, label style={font=\small},
]
\addplot[fill=sendcol, draw=sendcol] coordinates {(0.612,A) (0.534,L0) (0.494,S) (0.527,C) (0.333,L1) (0.330,E)};
\addplot[fill=allowcol, draw=allowcol] coordinates {(-0.024,A) (-0.030,L0) (0.116,S) (-0.001,C) (0.355,L1) (0.375,E)};
\legend{송금 검사, 허락 검사}
\end{axis}
\end{tikzpicture}
\caption{점수마다 두 검사와 잰 $\tau$. 위의 두 점수는 허락 지도를, 아래의 세 점수는 송금 지도를 걷고, 동전 랭크는 둘을 섞어 걸어요.}
\end{figure}

그래프의 점수는 위아래 두 무리로 나뉘어요.
위쪽의 엔도스랭크와 앞면 랭크는 허락 지도를 걸어요.
이 둘은 허락 검사와 닮았어요(0.375와 0.355).
아래쪽의 송금 랭크, 뒷면 랭크, 씨앗 랭크는 송금 지도를 걸어요.
이 셋은 송금 검사와 많이 닮았고(0.494\textasciitilde{}0.612), 허락 검사와는 거의 닮지 않았어요(씨앗 랭크만 0.116).
동전 랭크는 두 지도를 섞어 걷는데도 송금 쪽으로 기울었어요.
송금 화살표가 허락 화살표보다 아홉 배쯤 많아서, 쪽지가 대부분 송금 길로 다니기 때문이에요.

점수마다 자기가 걷는 지도의 기록과 닮은 거예요.
그런데 이건 어찌 보면 당연해요.
허락 검사는 허락 화살표를 센 것이고, 엔도스랭크는 바로 그 화살표 위에서 쪽지를 돌린 점수니까요.
그래서 이 결과로 “어느 점수가 더 좋다”고 말할 수는 없어요.
달리기 1등과 받아쓰기 1등이 다른 것처럼, 점수마다 재는 것이 다를 뿐이에요.

\section{우연은 아닐까?}

그런데 이런 차이가 우연히 나온 것일 수도 있잖아요?
그래서 과학자들은 \textbf{다시 뽑기}를 해 봐요.

지갑 5,521개가 적힌 쪽지를 상자에 넣고, 한 장 뽑아 적고 다시 넣기를 5,521번 해요.
그러면 어떤 지갑은 두 번 뽑히고 어떤 지갑은 안 뽑힌, 조금 다른 무리가 생겨요.
그 무리로 $\tau$를 다시 계산해요.
이것을 400번 되풀이하면 $\tau$가 400개 나와요.
그 가운데 가장 작은 쪽 2.5\%와 가장 큰 쪽 2.5\%를 빼고 남은 범위를 \textbf{95\% 범위}라고 해요.
“운이 조금 나쁘거나 조금 좋아도 대개 이 안에 들어온다”는 범위예요.

엔도스랭크와 송금 랭크의 $\tau$ 0.360은 95\% 범위가 0.342에서 0.379였어요.
범위가 좁고 1에서 멀리 떨어져 있으니, 두 점수가 다르다는 건 우연이 아니에요.
점수마다 검사와 잰 $\tau$에도 이런 범위를 붙였어요.
자기 지도의 검사와 잰 $\tau$는 범위가 모두 0에서 멀리 떨어져 있었어요.
이 다시 뽑기는 14장에서 다시 아주 중요해져요.

\begin{deeper}[진짜 이름과 숫자]
\begin{itemize}
\item 논문에서는 동점을 고려한 켄달의 $\tau_b$를 썼어요.
\item 다시 뽑기는 부트스트랩이라고 해요. 같은 지갑 표본으로 모든 점수를 함께 계산하는 “짝지은 부트스트랩”을 400번, 시드 42로 했어요.
\item 검증 지표는 여섯 가족으로 묶었어요: 송금, 허락, GMX 청산, 손실 위험, 시빌 안정성, GMX 거래 성과. 본문에서 주로 보는 것은 송금과 허락 가족이에요. 시빌 안정성은 모든 점수가 0.220에서 0.292 사이였고, GMX 쪽 세 가족과는 모든 점수가 약하게만 닮았어요(가장 높은 것도 0.236).
\item 송금 검사의 받은 송금 수와 허락 검사의 받은 허락 수는 4장에서 센 바로 그 수예요. 쪽지 놀이 점수는 이 수를 지도 전체로 퍼뜨린 것이라, 자기 지도의 검사와 닮는 것이 어느 정도 정해져 있어요.
\item 엔도스랭크의 0.375는 허락을 받은 131개 지갑을 기본 점수 위로 올려 세운 결과예요. 동점이 많아서 엔도스랭크가 허락 검사와 낼 수 있는 가장 높은 $\tau$가 0.38쯤이고, 동점이 없는 검사와는 0.57쯤이에요. 그래서 엔도스랭크의 두 숫자(0.375와 0.330)는 서로 견주지 않아요. 지도 밖 지갑을 화살표 없는 점으로 넣으면 허락 검사와는 0.990, 송금 검사와는 $-0.071$이 돼요.
\item 엔도스랭크와 송금 랭크 둘만 맞대어 잰 차이는 책 끝의 「엔도스랭크와 송금 랭크를 맞대어 보면」(\pageref{ch:er-awp-side}쪽)에 따로 모았어요. 두 점수는 재는 것이 달라서, 그 차이로 어느 쪽이 더 좋다고 말할 수 없기 때문이에요.
\item 엔도스랭크와 송금 랭크 점수는 이 검사들과 봄 채점(13장)의 결과가 나온 뒤에 매겼어요. 그래서 이 두 점수가 들어간 비교는 미리 정해 둔 것이 아니고, 범위와 함께 그대로 보고해요.
\end{itemize}
\end{deeper}

\begin{learned}
\begin{itemize}
\item 켄달의 $\tau$는 두 줄에서 맞는 쌍과 어긋난 쌍을 세어 두 순서가 얼마나 닮았는지 재요.
\item 엔도스랭크와 송금 랭크의 $\tau$는 0.360이에요. 두 점수는 지갑을 꽤 다르게 줄 세워요.
\item 점수마다 자기가 걷는 지도의 검사와 닮았어요. 검사가 그 지도의 화살표를 센 것이라 어느 정도는 당연한 결과예요.
\item 다시 뽑기를 400번 해서 차이가 우연이 아닌지 확인했어요.
\end{itemize}
\end{learned}
